\section{Experiments}

\subsection{Protocol and Metrics}

We organize the main evidence around four questions: whether CLE induces binding-aligned behavior; whether formation is predominantly local-first; whether PEW+BER suppresses DSA beyond generic ERM/group/tail-risk alternatives; and whether suppression persists under controlled changes of the binding map or communication base.

DSA is the primary shortcut metric. We also report operator-grid accuracy and frozen task utility metrics: client-average accuracy (Avg), minimum client accuracy (Worst), worst valid class--operator-cell accuracy (WCCA), and mean within-class accuracy range across operators (CFG, lower is better). Absolute values across different protocols are not compared directly; attribution relies on matched within-setting contrasts. This distinction is especially important because the original-map intervention experiment, the held-out map2 control experiment, and the FedDF-fidelity experiment intentionally use different controlled recipes. We therefore use each experiment to answer a specific scientific question rather than ranking all checkpoints in a single global leaderboard. The held-out map2 table is the primary comparison for alternative risk objectives; map1 tests dependence on one particular binding assignment; and FedDF-fidelity tests whether the shortcut-suppression effect is tied to the primary HFL communication stack. This organization prevents protocol changes from being misinterpreted as treatment effects.

\begin{table*}[t]
\centering
\small
\caption{Controlled protocols and information boundary. Operator metadata and binding maps are sealed from deployable training and used only for controlled construction or post-training evaluation.}
\label{tab:protocol-boundary}
\resizebox{\textwidth}{!}{%
\begin{tabular}{lllllll}
\toprule
Protocol & Private task & Public PEW carriers & Architectures & Rounds & Train-time environment information & Evaluation-only information \\
\midrule
Primary CLE / M1--M2 & CIFAR--10 & CIFAR--100 & R10/R12/ShuffleNet/MobileNetV2 & 12 or 40 & PEW family proxy & operator, binding, paired grid \\
Second task / M3 & CIFAR--100 & CIFAR--10 & same heterogeneous set & 40 & newly trained PEW family proxy & operator, binding, paired grid \\
Taxonomy stress / S1 & CIFAR--10 & CIFAR--100 minus motion blur & same heterogeneous set & 40 & leave-one-operator-out PEW proxy & held-out operator, binding, paired grid \\
HFL context / S2 & CIFAR--10 & protocol-specific public data & same heterogeneous set & 40 & method-specific, no BER & operator, binding, paired grid \\
\bottomrule
\end{tabular}%
}
\end{table*}


\begin{figure*}[t]
\centering
\includegraphics[width=0.98\textwidth]{figures/fig1_cle_hfl_evidence_chain.pdf}
\caption{CLE-HFL evidence chain and information boundary. Training-time PEW+BER uses only a coarse family proxy, whereas operator metadata and the binding map are reserved for sealed, reporting-only DSA evaluation.}
\label{fig:evidence-chain}
\end{figure*}

\subsection{Held-Out Matched Controls}

A five-arm 12-round screen was used only for design and method selection; it is not treated as final scientific evidence. In particular, JTT remains a screening-only comparison and we do not claim a Formal empirical victory over JTT.

The final held-out map2 comparison uses a frozen 40-round protocol with ERM, protocol-matched CVaR-DRO, PEW+GroupDRO, and PEW+BER. The four arms share the same partition, training/evaluation seed, initial states, private batch trajectory, and strict HFL communication protocol. PEW+GroupDRO and PEW+BER receive the same PEW grouping information.

\begin{table*}[t]
\centering
\small
\caption{Held-out map2 matched method comparison. Numeric entries are completed training-seed-0 Formal results; M1 will add seeds 1/2. JTT remains screening-only until O1 Formal is completed.}
\label{tab:map2-methods}
\resizebox{\textwidth}{!}{%
\begin{tabular}{lrrrrrr}
\toprule
Arm & DSA $\downarrow$ & Operator-grid Acc. $\uparrow$ & Last-10 Avg $\uparrow$ & Last-10 Worst $\uparrow$ & WCCA $\uparrow$ & CFG $\downarrow$ \\
\midrule
ERM & 0.260308 & 20.3067 & 19.8578 & 16.3007 & 0.000 & 32.8325 \\
CVaR-DRO & 0.088982 & 20.4567 & 18.6700 & 14.0380 & 0.125 & 27.7775 \\
JTT Formal & \pendingcell & \pendingcell & \pendingcell & \pendingcell & \pendingcell & \pendingcell \\
PEW+GroupDRO & 0.214575 & 20.7883 & 20.9562 & 17.8140 & 0.025 & 31.5275 \\
\textbf{PEW+BER} & \textbf{0.077741} & \textbf{23.7433} & \textbf{24.4518} & \textbf{20.4727} & \textbf{2.650} & \textbf{22.8025} \\
\bottomrule
\end{tabular}%
}
\ifdraftskeleton
\par\smallskip
\evidenceneeded{M1 Formal training seeds 1/2 for mean$\pm$std; O1 JTT Formal only if the final paper makes an empirical JTT comparison.}
\fi
\end{table*}


For training seed 0, source-paired bootstrap differences are

\[
\mathrm{DSA(ERM)}-\mathrm{DSA(BER)}=0.182567
\]

with 95\% interval \([0.180596,0.184407]\),

\[
\mathrm{DSA(GroupDRO)}-\mathrm{DSA(BER)}=0.136834
\]

with \([0.134952,0.138706]\), and

\[
\mathrm{DSA(BER)}-\mathrm{DSA(CVaR)}=-0.011240
\]

with \([-0.012355,-0.010092]\).

Across training seeds 0/1/2, DSA(ERM)$-$DSA(BER) is (0.179060\pm0.008594), and DSA(GroupDRO)$-$DSA(BER) is (0.133947\pm0.011518); both contrasts have the favorable direction in all three seeds. DSA(BER)$-$DSA(CVaR) is (-0.015184\pm0.015411), while BER improves operator-grid accuracy over CVaR by (5.9567\pm2.3151) percentage points. All pre-registered multi-seed gates pass, yielding \texttt{GO\_MAP2\_TRAINING\_SEED\_STABILITY}.

The matched PEW+GroupDRO result shows that shared PEW grouping alone is insufficient to reproduce BER's lower pooled DSA under this protocol. Relative to CVaR-DRO, BER has lower pooled DSA and substantially higher pooled task utility on held-out map2. The conclusion is not uniform at client level: in the three-seed means, CVaR-DRO still has lower DSA on clients c0 and c1, while BER has lower DSA on c2 and c3. We therefore state only that \textbf{PEW+BER achieves a better pooled shortcut--utility trade-off than CVaR-DRO on held-out map2, without uniformly dominating CVaR-DRO in per-client DSA.}

\subsection{Controlled Replication and Communication-Base Boundary}

We next summarize matched PEW+BER comparisons across completed binding-map and communication-base settings, while reserving the second-private-task row for M3.

\begin{table*}[t]
\centering
\small
\caption{Cross-binding-map, communication-base, and private-task summary. Completed rows are within-setting matched comparisons and must not be read as a global leaderboard.}
\label{tab:cross-setting-dataset}
\resizebox{\textwidth}{!}{%
\begin{tabular}{lllrrrr}
\toprule
Private task & Binding / protocol & HFL base & Base DSA & PEW+BER DSA & Relative reduction & Avg delta \\
\midrule
CIFAR--10 & original map & Asymmetric HFL & 0.119644 & 0.041252 & 65.52\% & +1.7323 \\
CIFAR--10 & map1 & Asymmetric HFL & 0.113761 & 0.049531 & 56.46\% & +0.7040 \\
CIFAR--10 & held-out map2 (seed 0) & Asymmetric HFL & 0.260308 & 0.077741 & 70.13\% & +4.5940 \\
CIFAR--10 & original map & FedDF-fidelity & 0.136989 & 0.029012 & 78.82\% & -0.6660 \\
CIFAR--100 & frozen second-task map & Asymmetric HFL & \pendingcell & \pendingcell & \pendingcell & \pendingcell \\
\bottomrule
\end{tabular}%
}
\ifdraftskeleton
\par\smallskip
\evidenceneeded{M1 seeds 1/2 for the held-out map2 summary and M3 benchmark/Formal for the CIFAR--100 row.}
\fi
\end{table*}


Matched utility deltas are not uniformly positive. On the original map they are Avg \(+1.7323\) and Worst \(-0.3760\); on map1 they are Avg \(+0.7040\) and Worst \(+1.0760\). Under FedDF-fidelity, operator-grid accuracy changes by \(-0.0400\) points and Avg by \(-0.6660\), and the frozen utility gate fails. We therefore use Table~\ref{tab:cross-setting-dataset} only to summarize within-setting DSA suppression, not to claim uniform utility improvement or a global leaderboard.

On map1, only the client-specific class--operator binding map is changed while the partition, evaluation seed, training seed, initial weights, public data, evaluation grid, and frozen PEW remain fixed. All four clients reduce DSA, and pooled operator-grid accuracy changes from \(20.9450\%\) to \(21.4583\%\). This supports replication across two controlled binding directions at seed 0; it does not establish cross-partition, cross-seed, cross-dataset, or real-world generalization.

The second communication-base experiment uses native single-view CE with a protocol-matched \textbf{FedDF-fidelity} adapter. Base and BER arms share initialization, private batches, public distillation, and training budget. DSA changes from \(0.136989\) to \(0.029012\), but utility is mixed: operator-grid accuracy changes by \(-0.0400\) points, last-five Avg by \(-0.6660\), and the frozen overall utility gate fails; both arms also fall below the pre-registered 20\% operator-grid learning floor. The safe interpretation is therefore: \textbf{the DSA-suppression effect is also observed under a second protocol-matched HFL communication base, although its utility gate fails.} This is not evidence for a lossless universal plugin or consistent cross-base utility improvement.

\subsection{HFL Context Baselines}

The final paper separates two comparison axes. Axis I, Table~\ref{tab:map2-methods}, compares generic shortcut and group-risk objectives under one matched HFL protocol. Axis II, Table~\ref{tab:hfl-context}, asks whether CLE behavior is visible across representative heterogeneous-FL mechanism families: no collaboration, public-logit exchange, server distillation, prototype communication, personalized transfer, cross-correlation collaboration, label-noise robustness, and corruption-robust HFL. Fidelity labels describe the implementation scope and do not imply exact reproduction of every official training budget. This protocol-matched 40-round table is distinct from the 40-local-pretraining-plus-40-collaboration recipe used by AugHFL/RAHFL in their original studies.

\ifdraftskeleton
\begin{table*}[t]
\centering
\small
\caption{Planned protocol-matched HFL context table. All metric cells remain pending until S2 benchmark and Formal complete; this table is not an official-recipe leaderboard.}
\label{tab:hfl-context}
\begin{tabular}{llllrrr}
\toprule
Method & Model heterogeneity & Local robustness & Communication & DSA $\downarrow$ & Operator Acc. $\uparrow$ & Avg $\uparrow$ \\
\midrule
Local/ERM & yes & none & none & \pendingcell & \pendingcell & \pendingcell \\
FedDF-fidelity & yes & none & public distillation & \pendingcell & \pendingcell & \pendingcell \\
KT-pFL-fidelity & yes & none & knowledge transfer & \pendingcell & \pendingcell & \pendingcell \\
FCCL protocol adapter & yes & none & logit collaboration & \pendingcell & \pendingcell & \pendingcell \\
AugHFL-fidelity & yes & AugMix/JSD & heterogeneous collaboration & \pendingcell & \pendingcell & \pendingcell \\
RAHFL anchor & yes & AugMix/JSD/DCL & asymmetric collaboration & \pendingcell & \pendingcell & \pendingcell \\
\bottomrule
\end{tabular}
\par\smallskip
\evidenceneeded{S2 benchmark-first, followed by separately authorized 40-round Formal.}
\end{table*}
\fi


\begin{figure*}[t]
\centering
\includegraphics[width=0.98\textwidth]{figures/fig2_dsa_proxy_boundary.pdf}
\caption{Paired directional diagnosis and proxy-identifiability boundary. DSA compares bound-class probability mass across operators for the same source; identical observable label--proxy evidence can correspond to latent worlds with different true dependence.}
\label{fig:dsa-proxy}
\end{figure*}

\subsection{Environment Correspondence and Granularity}

A non-deployable Oracle audit compares true family grouping, true operator grouping, and a random operator correspondence. Correct environment correspondence materially affects DSA in this controlled audit: Random operator DSA exceeds Oracle operator DSA by \(0.052112\), with source-bootstrap 95\% interval \([0.051231,0.052983]\). The Oracle family--operator DSA gap is only \(0.001449\), with \([0.001086,0.001790]\). This does \textbf{not} establish that operator-level environment models are universally unhelpful. It only means that the present Oracle granularity experiment does not provide sufficient evidence to justify developing an operator-level PEW for this paper. The complete Oracle table is moved to Appendix C.

\subsection{Remaining Evidence Placeholders}

The remaining table cells are deliberately marked pending. M1 is complete. M2 is an implemented Formal repetition; M3 is an implemented second-private-task protocol that still requires benchmark-first execution; S1 is a bounded leave-one-operator-out taxonomy stress; S2 is a protocol-matched HFL context table; and O1 is a conditional JTT Formal. None of M2/M3/S1/S2/O1 should be described as launched, in progress, or successful until returned artifacts are audited. The paper does not add a new loss, communication mechanism, hierarchical PEW, or operator-level PEW in response to these gaps.
